\documentclass[11pt]{article}

\usepackage[utf8]{inputenc}
\usepackage[T1]{fontenc}
\usepackage{hyperref}
\usepackage{booktabs}
\usepackage{listings}
\usepackage{geometry}
\usepackage{microtype}
\usepackage{amsmath}

\geometry{margin=1in}

\setlength{\parindent}{0pt}
\setlength{\parskip}{6pt}

\hypersetup{
    colorlinks=true,
    linkcolor=blue,
    urlcolor=blue,
    citecolor=blue
}

\lstset{
    basicstyle=\ttfamily\small,
    breaklines=true,
    frame=single,
    backgroundcolor=\color{gray!10},
}
\usepackage{xcolor}

\title{Named Entity Recognition for Ancient Egyptian Transliteration}

\author{
    Juhi Jadhav \\
    Independent Researcher \\
    \texttt{juhij3.me@gmail.com}
}

\date{}

\begin{document}

\maketitle

\begin{abstract}
Named entity recognition (NER) for ancient Egyptian texts has been limited to work on modern translations, leaving the source language largely unaddressed by computational methods. This paper presents the first NER model trained directly on ancient Egyptian Leiden Unified Transliteration, identifying divine and human entities without reliance on English translation. Using the Thesaurus Linguae Aegyptiae (TLA) corpus of 12,773 Middle Egyptian sentences, training examples are automatically generated by leveraging existing UPOS part-of-speech annotations, eliminating the need for manual entity labeling. A language-agnostic spaCy NER model trained on 7,059 annotated sentences achieves 95.6\% F1 across DEITY and PERSON entity classes. The model is demonstrated on Book of the Dead Spell 125 and implications for automated entity extraction pipelines connecting transliteration text to structured demonological knowledge bases such as DemonThings/DemonBase are discussed.
\end{abstract}

\section{Introduction}

Ancient Egyptian religious texts, including the Pyramid Texts, Coffin Texts, and the Book of the Dead, contain dense networks of divine and demonic entities. Identifying these entities computationally is a prerequisite for large-scale analysis of how deities, demons, and ritual agents are distributed across spells, time periods, and manuscript traditions.

Existing computational work on ancient Egyptian texts has primarily targeted two problems: hieroglyph sign classification \cite{barucci2021,franken2013} and machine translation from hieroglyphic encodings to modern languages \cite{hieroglyphic2024}. NER work in the broader Digital Humanities community has addressed Latin \cite{romanello2023}, Greek, and Chinese historical texts, but ancient Egyptian transliteration has received no attention.

The dominant reason for this gap is the absence of annotated training data. Manually annotating entity spans in Egyptian transliteration requires rare specialist expertise combining Egyptological knowledge and familiarity with the Leiden Unified Transliteration (LUT) system, a standardized ASCII encoding of Egyptian phonological values used by professional Egyptologists worldwide.

This paper addresses the gap through a simple but effective observation: the Thesaurus Linguae Aegyptiae (TLA), the most comprehensive digital corpus of ancient Egyptian texts, already contains expert part-of-speech annotations including the PROPN (proper noun) tag. These annotations were produced by trained Egyptologists as part of the TLA lemmatization project and provide a high-quality training signal for NER without requiring additional manual annotation.

The contributions of this paper are:

\begin{enumerate}
    \item The first NER model for ancient Egyptian Leiden Unified Transliteration
    \item A method for automatically generating NER training data from existing UPOS annotations in the TLA corpus
    \item A publicly available trained model and interactive demo
    \item A demonstration on Book of the Dead Spell 125 with analysis of model behavior on real manuscript text
\end{enumerate}

\section{Related Work}

\subsection{Computational Egyptology}

Hieroglyph recognition has been the primary focus of computational work on ancient Egyptian. Franken \& van Gemert \cite{franken2013} introduced the GlyphReader dataset and a retrieval-based classifier for individual hieroglyph signs. Barucci et al.\ \cite{barucci2021} extended this with deep learning approaches achieving strong classification accuracy across Gardiner sign classes. More recently, the OCR-PT-CT project \cite{ocrptct2022} applied metric learning to semi-automatic transcription of Pyramid Texts and Coffin Texts, achieving 97.7\% accuracy on sign detection.

At the language level, the Hieroglyphic Transformer \cite{hieroglyphic2024} built on the M2M-100 framework to enable end-to-end translation from hieroglyphic encodings to English, achieving 86.3\% accuracy on authentic inscriptions. The BBAW Egyptian dataset provides parallel hieroglyphic encodings, transliterations, and translations for 35,503 sentences.

None of this work addresses entity extraction from transliteration text.

\subsection{NER for Ancient and Historical Languages}

NER for classical languages has received growing attention. Romanello et al.\ \cite{romanello2023} developed NER models for classical Latin and Greek, focusing on person and place names in historical texts. Work on ancient Chinese \cite{chisiec2024} demonstrated that BERT-based approaches can achieve strong performance on historical NER with domain-specific pretraining. Recent work has explored few-shot and zero-shot NER for low-resource ancient languages using large language models \cite{ml4al2024}.

Egyptian transliteration presents unique challenges absent in these languages: non-standard character sets (\textit{\=a}, \textit{i\kern-1pt{}}, \d{h}, \d{H}, \d{H}\kern-1pt{}, \d{t}, \d{d}), compound entity names (e.g., \texttt{nbt-\d{h}w.t} for Nephthys), grammatical affixes attached directly to name tokens (e.g., \texttt{ws\textit{i\kern-1pt{}}r=f}), and systematic orthographic variation across time periods and scribal traditions.

\subsection{Egyptian Demonology Databases}

The DemonThings project \cite{scalf} maintains DemonBase, a relational database of 4,043 liminal entities from ancient Egyptian sources including coffins, wands, headrests, and complete Book of the Dead manuscripts. Sartini \& Lucarelli \cite{sartini2024} published a semantic knowledge graph representation of this database, providing the first structured ontology for Egyptian demonology. These databases represent the downstream application context for the NER pipeline presented here.

\section{Dataset}

\subsection{Thesaurus Linguae Aegyptiae}

The Thesaurus Linguae Aegyptiae (TLA) \cite{tla} is maintained by the Berlin-Brandenburg Academy of Sciences and Humanities and represents the most comprehensive digital corpus of ancient Egyptian texts. The corpus covers texts from the Old Kingdom through the Roman Period (ca.\ 2700 BCE to 400 CE) across hieroglyphic, hieratic, and Demotic scripts.

The dataset used in this paper is \texttt{thesaurus-linguae-aegyptiae/tla-Earlier\_Egyptian\_original-v18-premium}, available on HuggingFace under CC BY-SA 4.0. It contains 12,773 sentences of Middle and Earlier Egyptian with the following fields:

\begin{itemize}
    \item \textbf{hieroglyphs}: Unicode hieroglyphic encoding
    \item \textbf{transliteration}: Leiden Unified Transliteration
    \item \textbf{lemmatization}: TLA lemma IDs with lemma forms
    \item \textbf{UPOS}: Universal part-of-speech tags
    \item \textbf{glossing}: Grammatical glosses
    \item \textbf{translation}: German translation
    \item \textbf{dateNotBefore / dateNotAfter}: Dating in years BCE/CE
\end{itemize}

The corpus spans the period ca.\ 2180--1539 BCE (First Intermediate Period through Second Intermediate Period), corresponding to the primary period of Coffin Text and early Book of the Dead composition.

\subsection{Transliteration Encoding}

Leiden Unified Transliteration uses a standardized set of special characters to represent Egyptian phonemes. Key characters include \d{a} (aleph), \textit{\i} (yod), \d{h} (emphatic h), \d{H} (velar fricative), \b{H} (palatal fricative), \d{t} (palatal t), and \d{d} (emphatic d). Entity names in this encoding differ substantially from their common English forms: Osiris appears as \texttt{ws\textit{i}r}, Anubis as \texttt{\textit{i}np.w}, Horus as \texttt{\d{h}r.w}, and Seth as \texttt{st\v{s}}.

\section{Method}

\subsection{Automatic Training Data Generation}

Manual annotation of entity spans in Egyptian transliteration requires expert Egyptological knowledge and is prohibitively expensive at scale. This paper exploits the existing UPOS annotations in the TLA corpus as an automatic training signal.

The TLA corpus was annotated by trained Egyptologists and contains UPOS tags for all tokens. The PROPN (proper noun) tag identifies person names, deity names, place names, and other proper entities. Of 12,773 sentences, 9,187 PROPN tokens were identified across the corpus.

For each sentence, PROPN-tagged tokens are extracted with their character-level offsets in the transliteration string. A coarse entity type distinction is applied: tokens matching a seed list of 32 known deity names are labeled DEITY; all other PROPN tokens are labeled PERSON. This yields 7,059 training sentences with at least one labeled entity.

The seed deity list was compiled from standard Egyptological references and includes major deities attested across the corpus period: Osiris (\texttt{ws\textit{i}r}), Horus (\texttt{\d{h}r.w}), Ra (\texttt{r\d{a}w}), Anubis (\texttt{\textit{i}np.w}), Seth (\texttt{st\v{s}}), Geb (\texttt{gbb}), Nut (\texttt{nw.t}), Isis (\texttt{\textit{i}s.t}), Thoth (\texttt{\d{d}\d{h}wt\textit{i}}), Ptah (\texttt{pt\d{h}}), Sekhmet (\texttt{s\d{H}m.t}), Amun (\texttt{\textit{i}mn}), Hathor (\texttt{\d{h}wt-\d{h}r}), and others.

\subsection{Model Training}

A language-agnostic blank spaCy model (\texttt{spacy.blank("xx")}) is used as the base, avoiding any assumptions about tokenization or morphology that would be inappropriate for Egyptian transliteration. An NER pipeline component is added and trained from scratch.

The dataset is split 80/20 into training (5,647 sentences) and validation (1,412 sentences). The model is trained for 20 epochs with a batch size of 32 using spaCy's default Adam optimizer. The best model by validation F1 is saved after each epoch.

\section{Experiments and Results}

\subsection{Training Performance}

The model converges rapidly, achieving 92.1\% F1 on the validation set after the first epoch. Performance improves steadily across 20 epochs, with the best model achieving 95.6\% F1 (Precision: 96.0\%, Recall: 95.3\%) at epoch 17.

\begin{table}[h]
\centering
\begin{tabular}{@{}rrrrrr@{}}
\toprule
Epoch & Loss & Precision & Recall & F1 \\
\midrule
1  & 6695.5 & 93.8\% & 90.6\% & 92.1\% \\
5  & 285.9  & 93.8\% & 93.7\% & 93.8\% \\
10 & 176.6  & 96.2\% & 94.2\% & 95.2\% \\
17 & 100.2  & 96.0\% & 95.3\% & \textbf{95.6\%} \\
20 & 100.9  & 96.5\% & 94.4\% & 95.5\% \\
\bottomrule
\end{tabular}
\caption{Training performance across epochs. Best model saved at epoch 17.}
\label{tab:training}
\end{table}

\subsection{Evaluation on Book of the Dead Spell 125}

Spell 125 of the Book of the Dead, the Negative Confession, contains one of the densest concentrations of named divine entities in Egyptian funerary literature. The deceased addresses 42 divine assessors by name. The following passage was used as an out-of-domain test:

\begin{lstlisting}
ws'r `nh-ms-n'w m3`-hrw dd.'  n 'np.w nb t3-dsr hr.w nb p.t
r`w ntr `3 's.t mw.t ntr.t `3.t nbt p.t gbb nb t3.wj sts ntr
hq3.w dhwtj nb hmnw nbt-hw.t nbt 'mnt.t pth nb m3`.t 'mn.w
nb nswt t3.wj shm.t 'r.t-r`w
\end{lstlisting}

The model correctly identified: Osiris (\texttt{ws\textit{i}r}), Anubis (\texttt{\textit{i}np.w}), Horus (\texttt{\d{h}r.w}), Ra (\texttt{r\d{a}w}), Geb (\texttt{gbb}), Seth (\texttt{st\v{s}}), Thoth (\texttt{\d{d}\d{h}wt\textit{i}}), Ptah (\texttt{pt\d{h}}), Sekhmet (\texttt{s\d{H}m.t}) as DEITY, and \texttt{\d{a}n\d{h}-ms-n\textit{i}w} as PERSON (correctly identifying the deceased's name).

\subsection{Error Analysis}

Three systematic error types were observed:

\textbf{Compound deity names:} \texttt{nbt-\d{h}w.t} (Nephthys) was labeled PERSON rather than DEITY. Compound names formed with common nouns (\texttt{nbt} = lady, \texttt{\d{h}wt} = house) are not in the seed deity list and the model has not learned to generalize across compound forms.

\textbf{Divine epithets:} \texttt{\textit{i}r.t-r\d{a}w} (Eye of Ra) was labeled PERSON. Epithets functioning as entity names represent a genuinely ambiguous category, as they are not deities themselves but divine manifestations.

\textbf{Place names labeled PERSON:} \texttt{t\d{a}.wj} (the Two Lands, i.e., Egypt) was labeled PERSON. The current two-class taxonomy lacks a LOCATION class, causing place names to fall into PERSON by default.

These error types suggest clear directions for extension: expanding the seed deity list with compound forms, adding a LOCATION class, and incorporating a DIVINE\_EPITHET subclass.

\section{Discussion}

\subsection{Automatic Annotation Quality}

The use of TLA UPOS annotations as automatic training labels introduces a specific limitation: the DEITY/PERSON distinction is applied heuristically via a seed list rather than through expert annotation. The quality of the DEITY labels is therefore bounded by the completeness of the seed list. For the major deities of the Middle Kingdom period covered by this corpus, the seed list achieves good coverage. For minor deities, demon names, and composite divine entities, coverage is incomplete.

This limitation is precisely the research gap that motivates future collaboration with Egyptological experts and databases such as DemonBase. The DemonBase corpus of 4,043 entities, if made available for training, would allow the seed list to be replaced with a comprehensive authority file and the entity taxonomy to be extended to include DEMON, LOCATION, and RITUAL classes.

\subsection{Connection to DemonThings Pipeline}

The long-term vision for this work is an automated pipeline connecting manuscript images to structured demonological knowledge:

\begin{lstlisting}
Papyrus image
    -> Hieroglyph sign detection (OpenCV contour analysis)
    -> Sign classification (ResNet-50, 91.8% accuracy)
    -> Transliteration sequence
    -> NER (this model, 95.6% F1)
    -> Entity linking to DemonBase
    -> Knowledge graph population
\end{lstlisting}

Each component of this pipeline has been demonstrated independently. The NER model presented here provides the critical link between transliteration text and structured entity databases.

\section{Conclusion}

This paper presented the first NER model for ancient Egyptian Leiden Unified Transliteration, achieving 95.6\% F1 on DEITY and PERSON entity classes without requiring manual annotation. The method leverages existing UPOS annotations in the TLA corpus as an automatic training signal, making the approach reproducible and extensible to other annotated ancient language corpora.

The model is publicly available at \url{https://github.com/Juhij2/egyptian-transliteration-ner} with an interactive demo at \url{https://juhi-egyptian-transliteration-ner.streamlit.app/}. This preprint is available at \url{https://doi.org/10.5281/zenodo.19800720}.

Future work includes extending the entity taxonomy to DEMON, LOCATION, and RITUAL classes; improving compound name recognition; and integrating the model into a full manuscript-to-knowledge-graph pipeline in collaboration with the DemonThings project at the Oriental Institute, University of Chicago.

\bibliographystyle{plain}
\bibliography{references}

\end{document}
